Работу выполнил И.~И.~Иванов под руководством Петрова~П.~П.

Слово перед инициалами привязывают, только если фамилия стоит перед ними:
отчёт подготовил И.~И.~Иванов, проверил Петров~П.~П.

Одиночная прописная буква в конце предложения инициалом не считается:
протоколы приведены в приложении А. В работе они не разбираются.

Сокращение из нескольких букв инициалом не считается: см. рис. 1.

Инициалы внутри буквальной вставки не трогают: \verb|И. И. Иванов|.
